%!TEX root = ../main.tex
\section{System Overview}
\label{sec:overview}

%!TEX root = ../../main.tex
% Placeholder: framework overview figure (two-column span), to be hand-drawn.
\begin{figure*}[t]
    \centering
    \setlength{\fboxrule}{0.8pt}%
    \textcolor{gray!55}{\framebox[\dimexpr\textwidth-2\fboxsep-2\fboxrule\relax]{%
      \begin{minipage}[c][0.24\textwidth][c]{0.95\textwidth}
        \centering
        \textcolor{gray!65}{\normalsize\itshape [Placeholder] Framework overview of \sys{}}\\[0.6em]
        \textcolor{gray!55}{\small catalog $+$ planner $+$ operators
          $\;\longrightarrow\;$ second-order memory $\;\longrightarrow\;$ cost-aware executor}
      \end{minipage}}}
    \caption{Overview of \sys{}. A foundational framework (catalog, planner,
      operators) is augmented by two optimization layers:
      the Second-Order Memory, which closes a workload feedback loop,
      and the Cost-Aware Executor, which combines pattern-aware budgeting with
      agreement-based early stopping.}
    \Description{Placeholder box for the system architecture diagram of
      MemCatalog, showing the catalog, planner, and operators
      as the base framework, with the second-order memory and the cost-aware
      executor layered on top.}
    \label{fig:framework_overview}
\end{figure*}


This section traces one query end to end, so that the division of labor among the three mechanisms is visible before any of them is specified in detail. Figure~\ref{fig:framework_overview} shows the resulting architecture, in which a framework supplies the plan and two optimization layers refine what that plan reads and what it costs.

\stitle{Life Cycle of a Query.}
Given a query $q$, \sys{} evaluates Eq.~(\ref{eq:sys}) in four steps. It first \emph{plans} the query, mapping $q$ to one of \point, \agg, or \cmp{} from surface form alone (Section~\ref{subsec:planner}). It then \emph{grounds} the selected plan, resolving the catalog indexes $\mathcal{I}_e$ and $\mathcal{I}_a$ that the chosen operator will read, with the second-order memory reweighting those indexes by the per-user priority $\pi(e)$ so that entities this user has already confirmed appear first (Section~\ref{subsec:feedback}). It then \emph{executes}, issuing LLM calls under the operator's fixed template while the executor decides how many calls the plan class warrants and stops early once two samples agree (Sections~\ref{subsec:budget} and~\ref{subsec:shortcircuit}). It finally \emph{records} the outcome, writing the entities that appeared in the accepted answer into the workload log $\mathcal{W}_u$ so that the next query starts from a better-ordered catalog (Section~\ref{subsec:log}).

\stitle{Interlocking of the Three Mechanisms.}
The ordering above is not a pipeline of independent stages, since each mechanism consumes something another produces. The executor cannot price a query until the planner has named its class. The planner cannot read a useful candidate set until the second-order memory has ranked the catalog. The second-order memory cannot update until the executor has committed to an answer worth recording. Removing any one therefore degrades the other two, which Section~\ref{sec:experiments} confirms by ablation. The remainder of the paper takes the three in turn.
